% STATUS: reviewed
\chapter{Konfigūracija ir derinamos konstantos}\label{app:config}

Šis priedas --- žinynas apie kiekvieną reikšmę, kurią gali nustatyti diegimas, ir kiekvieną kodo
konstantą, nuo kurios priklauso, ką modeliai gauna, ką jie daro ir kiek tai kainuoja. Pirmiausia
parodoma, iš kur reikšmė atsiranda (\cref{sec:config:where}), tada išvardijami API ir asistento
paslaugos aplinkos kintamieji su numatytosiomis reikšmėmis, kurias jiems suteikia kodas
(\cref{sec:config:api,sec:config:assistant}), ir įvardijami kitų konteinerių kintamieji
(\cref{sec:config:other}). Ilgiausia dalis --- derinamų konstantų lentelė, suskirstyta pagal
sritis (\cref{sec:config:constants}). Pabaigoje aprašoma, kaip keičiama konstanta ir kurios
instrukcijos bei testai nuo jos priklauso (\cref{sec:config:change}). Kaip su šiais parametrais
eksploatuoti sistemą, aprašyta \cref{ch:ops}; kodėl kiekviena konstanta turi būtent tokią
reikšmę, paaiškinta jos srities skyriuje. Visos čia pateiktos reikšmės perskaitytos iš kodo
2026~m. lapkričio 13~d. Slapti duomenys įvardijami tik pagal paskirtį, niekada --- su reikšme.


\section{Iš kur atsiranda reikšmė}\label{sec:config:where}

Parametras veikiančią paslaugą pasiekia vienu iš kelių būdų (\cref{fig:config:where}). Paties
operatoriaus pasirinkimai ir slapti duomenys yra serverio failo \file{.env} eilutės; jį Docker
Compose perskaito kurdamas konteinerį ir perduoda per to konteinerio \code{environment} bloką
faile \file{docker-compose.yml}. Kai kurios reikšmės įrašytos tiesiai į \textit{compose} failą.
Kiekviena paslauga savo aplinką perskaito vieną kartą, kai paleidžiama: API --- faile
\file{settings.py} (o žiniatinklio serveris --- faile \file{django/gunicorn.conf.py}), asistento
paslauga --- faile \file{config.js}. Nenustatytas kintamasis gauna tame kode įrašytą numatytąją
reikšmę. Derinamos konstantos išvis nėra parametrai: tai vardai juos naudojančių modulių
pradžioje, ir jos keičiasi tik kartu su kodu. Galiausiai trys kiekvieno kalbos modelio
iškvietimo savybės kode neįrašytos: modelis, temperatūra ir išvesties žetonų riba yra aktyvios
instrukcijos versijos laukai, kuriuos administratorius redaguoja konsolėje (\cref{ch:prompts});
be to, administratorius gali paskyrai nustatyti atskirą dienos biudžetą, kuris tada pakeičia
\code{LLM_DAILY_BUDGET_USD}.

\begin{figure}[H]
\centering
\begin{tikzpicture}[x=1mm, y=1mm]
  \node[ext, text width=24mm] (env) at (0,0) {\file{.env} serveryje\\slapti duomenys, pasirinkimai};
  \node[svc, text width=29mm] (compose) at (36,0) {\file{docker-compose.yml}\\kiekvieno\\konteinerio aplinka};
  \node[svc, text width=28mm] (api) at (80,9) {API\\\file{settings.py}};
  \node[svc, text width=28mm] (asst) at (80,-9) {asistento paslauga\\\file{config.js}};
  \node[ai, text width=18mm] (call) at (122,0) {modelio iškvietimas};
  \begin{scope}[on background layer]
    \node[lane, fit=(api)(asst)] (code) {};
  \end{scope}
  \node[lanelabel, anchor=south] at (code.north) {numatytosios reikšmės ir konstantos kode};
  \node[store, text width=24mm] (pv) at (122,-30) {instrukcijų versijos\\modelis,\\temperatūra,\\išvesties riba};
  \node[store, text width=24mm] (acc) at (80,-34) {paskyros\\atskiras dienos biudžetas};
  \draw[flow] (env) -- (compose);
  \draw[flow] (compose.east) -- ++(4,0) |- (api.west);
  \draw[flow] (compose.east) -- ++(4,0) |- (asst.west);
  \draw[flow] (api.east) -- ++(5,0) |- (call.west);
  \draw[flow] (asst.east) -- ++(5,0) |- (call.west);
  \draw[flow] (pv.north) -- (call.south);
  \draw[opt] (acc.north) -- (code.south);
  \node[note, anchor=north] at (36,-9) {nenustatyta: kodo numatytoji reikšmė};
\end{tikzpicture}
\caption{Iš kur atsiranda reikšmės, lemiančios modelio iškvietimą. Operatorius nustato
\file{.env} ir \textit{compose} failą; kiekviena paslauga savo aplinką perskaito vieną kartą,
kai paleidžiama, o trūkstamai reikšmei naudoja numatytąją iš savo kodo; derinamos konstantos
keičiasi tik su kodu; konsolėje redaguojama instrukcijos versija kiekvienam iškvietimui suteikia
modelį, temperatūrą ir išvesties ribą; atskiras paskyros biudžetas pakeičia viso diegimo
biudžetą.}\label{fig:config:where}
\end{figure}


\section{API aplinka}\label{sec:config:api}

\Cref{tab:config:api} išvardija visus aplinkos kintamuosius, kuriuos skaito API kodas. Stulpelis
\enquote{Compose} rodo, kaip kintamąjį perduoda \file{docker-compose.yml}: \enquote{.env}, kai
reikšmė imama iš \file{.env} (jei \textit{compose} failas turi savo atsarginę reikšmę, ji
nurodoma skliaustuose), \enquote{failas}, kai reikšmė įrašyta \textit{compose} faile, ir
brūkšnys, kai \textit{compose} failas kintamojo neperduoda ir galioja kodo numatytoji reikšmė.
Užduočių vykdytojas veikia kaip \texttt{manage.py run\_worker} tame pačiame konteineryje
(\cref{ch:arch}), todėl jo aplinka ta pati.

\begin{xltabular}{\textwidth}{@{}L{40mm}L{21mm}Y@{}}
\caption{API aplinkos kintamieji (\file{django/compliance/settings.py}, jei nenurodytas kitas
failas)}\label{tab:config:api}\\
\toprule
\hd{Kintamasis} & \hd{Compose} & \hd{Numatytoji reikšmė kode ir ką nustato} \\
\midrule
\endfirsthead
\toprule
\hd{Kintamasis} & \hd{Compose} & \hd{Numatytoji reikšmė kode ir ką nustato} \\
\midrule
\endhead
\bottomrule
\endlastfoot
\code{DJANGO_SECRET_KEY} & .env & Privalomas. Slaptas: Django pasirašymo raktas. Be jo API nepasileidžia, nebent įjungtas \code{DJANGO_DEBUG} --- tada vietoj jo naudojamas fiksuotas kūrimo raktas. API autentifikavimas jo nenaudoja (prieigos raktai). \\
\code{DJANGO_DEBUG} & -- & \code{False}. \code{True} --- tik vietiniam derinimui: Django derinimo puslapiai. \\
\code{DATABASE_URL} & failas & Privalomas. PostgreSQL ryšys (su pgvector) URL pavidalu. Jame yra duomenų bazės slaptažodis, todėl jis slaptas visur, kur nukopijuotas. \\
\code{DB_CONN_MAX_AGE} & -- & \code{60}; \code{0} su kūrimo serveriu (\code{DJANGO_RUNSERVER}). Kiek sekundžių duomenų bazės ryšys laikomas pakartotiniam naudojimui; prieš naudojant ryšys patikrinamas. Kūrimo serveris kiekvienai užklausai paleidžia naują giją, ir laikomas ryšys liktų be savo gijos. \\
\code{UPLOAD_DIR} & -- & \code{/data/uploads}. Saugomų failų aplankas (teisės aktų šaltinių failai ir įmonės dokumentai). \\
\code{UPLOAD_QUOTA_FILES} & .env (\code{1000}) & \code{1000}. Kiek dokumentų gali turėti viena paskyra. \\
\code{UPLOAD_QUOTA_BYTES} & .env (\code{2147483648}) & \code{2 * 1024 * 1024 * 1024}. Kiek baitų gali turėti viena paskyra (2~GiB). \\
\code{ASSISTANT_URL} & -- & \code{http://legalops-assistant:3000}. Asistento paslaugos adresas; kiekvienas modelio iškvietimas ir vektorizavimas eina į jos \code{/internal/llm/*} maršrutus. \\
\code{ASSISTANT_INTERNAL_SECRET} & .env & Tuščias. Slaptas: saugo vidinius maršrutus tarp API ir asistento paslaugos abiem kryptimis; kai tuščias, atmetamas kiekvienas vidinis iškvietimas. \\
\code{AI_EMBED_MODEL} & failas & \code{gateway:azure/text-embedding-3-large}; \textit{compose} faile įrašyta \code{azure/text-embedding-3-large}, tas pats modelis (identifikatorius be tiekėjo priešdėlio yra DI tinklų sąsajos identifikatorius). Vektorizavimo modelis, nurodomas kiekviename iškvietime, kuris išsaugo vektorių (teisės aktų fragmentų, pareigų, įmonės dokumentų), ir įrašomas vektorizavimo etape. Jį pakeitus, reikia perskaičiuoti visus vektorius (\cref{ch:ops}). \\
\code{LLM_DAILY_BUDGET_USD} & .env (\code{10.0}) & \code{10.0}. Kiek JAV dolerių viena paskyra gali išleisti modelių iškvietimams per UTC parą; \code{0} --- be ribos. Administratoriaus nustatytas atskiras paskyros biudžetas jį pakeičia; administratorius be atskiro biudžeto neribojamas (\code{budget_status} faile \file{assistant/usage.py}). \\
\code{LOGIN_RATE_PER_IP} & .env (\code{60}) & \code{60}. Prisijungimo bandymai iš vieno kliento adreso per minutę. Adresu laikoma antraštė \code{X-Real-IP}, kurią įėjimo taškas nustato iš gautos antraštės \code{X-Forwarded-For}; veikiančiame serveryje dauguma lankytojų ateina vienu adresu ir dalijasi šiuo langu (\cref{ch:security}). Paskyros riba, 10 bandymų per 600~s, fiksuota faile \file{users/api/auth_views.py} (\code{EMAIL_LIMIT}, \code{EMAIL_WINDOW_S}). \\
\code{GUNICORN_WORKERS}, \code{GUNICORN_THREADS} & .env (\code{3}, \code{8}) & \code{3}, \code{8}; skaitoma faile \file{django/gunicorn.conf.py}. Gamybinio žiniatinklio serverio procesai ir gijos. \\
\code{DJANGO_RUNSERVER} & failas (\code{True}) & Nenustatytas; skaitomas faile \file{django/Dockerfile}. \code{True} vietoj gunicorn paleidžia Django kūrimo serverį; ši \textit{compose} eilutė priklauso kūrimo režimui (\cref{ch:ops}). \\
\code{WORKER_HEARTBEAT_FILE} & -- & \code{/tmp/worker-heartbeat}; skaitomas faile \file{jobs/management/commands/run_worker.py}. Failas, kurį užduočių vykdytojo ciklas paliečia kiekviename prasisukime ir per ilgą užduotį, kad sveikatos patikra galėtų atskirti užstrigusį ciklą nuo gyvo; kol kas jo niekas neskaito. \\
\code{WORKER_PID_FILE} & -- & \code{/tmp/worker.pid}; skaitomas tame pačiame faile. Veikiančio užduočių vykdytojo ciklo proceso identifikatorius; naujas ciklas paprašo jame nurodyto ciklo sustoti ir perima darbą, todėl antram ciklui, kuris turėtų veikti greta, reikia savo failo. \\
\code{XDG_CACHE_HOME} & failas & Projekto kodas jo neskaito: \textit{compose} failas nustato jį į \code{/tmp/.cache}, kad PDF eksporto šriftų biblioteka turėtų rašomą podėlį, nes konteinerio šakninė failų sistema skirta tik skaityti. \\
\end{xltabular}

Ne skaičius ten, kur laukiama skaičiaus (pvz., \code{UPLOAD_QUOTA_FILES}), neleidžia įkelti
parametrų, todėl API nepasileidžia. Be kintamųjų, \file{settings.py} fiksuoja kelias
nekonfigūruojamas ribas: JSON užklausų turinys skaitomas atmintyje iki 8~MB
(\code{DATA_UPLOAD_MAX_MEMORY_SIZE}), o didesni nei 5~MB įkelti failai perkeliami į
\file{/tmp} (\code{FILE_UPLOAD_MAX_MEMORY_SIZE}); kiekviena užklausa vykdoma vienoje duomenų
bazės transakcijoje (\code{ATOMIC_REQUESTS}); laikas --- UTC. Gamybinio žiniatinklio serverio
laiko ribos (\code{timeout = 330}, \code{graceful_timeout = 25}, \code{keepalive = 65}
sekundės) yra faile \file{django/gunicorn.conf.py}. Kadangi serverio procesai dirba gijomis,
pirmoji suveikia tik tada, kai procesas nustoja siųsti gyvybės signalą, ir niekada dėl to, kad
viena užklausa lėta: ataskaitos eksportas trunka minutes, o instrukcijos bandomasis paleidimas
savo modelio laukia iki 930~s.


\section{Asistento paslaugos aplinka}\label{sec:config:assistant}

Asistento paslauga visą savo aplinką perskaito faile \file{config.js} vieną kartą, kai modulis
pirmą kartą importuojamas. \Cref{tab:config:assistant} išvardija visus kintamuosius, kuriuos ji
skaito. Iš jų \file{docker-compose.yml} perduoda tik pirmuosius šešis; kiti lieka su kodo
numatytąja reikšme, nebent operatorius juos įrašo į paslaugos \code{environment} bloką.

\begin{codebox}{Kaip \file{config.js} perskaito reikšmę}
Skaičius turi būti vien skaitmenys leistiname intervale (\code{intFromEnv}): nuo lentelėje
nurodyto minimumo iki 2\,147\,483\,647 (ilgiausias delsimas, kurį gali laikyti Node laikmatis),
o prievadui --- iki 65\,535. Bet kas kita --- žodis, neigiamas skaičius, \code{0} ten, kur
minimumas yra 1, --- pakeičiama numatytąja reikšme, o paleidžiant įrašomas įspėjimas, nes nulis
nutrauktų kiekvieną ėjimą dėl laiko ribos arba išjungtų užklausų dažnio ribą. Nenustatytas ar
tuščias kintamasis reiškia numatytąją reikšmę. Iš URL pašalinami galiniai pasvirieji brūkšniai;
raktai, slaptas raktas ir modelių identifikatoriai paliekami baitas į baitą, o vien iš tarpų
sudarytas raktas laikomas nenustatytu.
\end{codebox}

\begin{xltabular}{\textwidth}{@{}L{38mm}L{40mm}Y@{}}
\caption{Asistento paslaugos aplinkos kintamieji (\file{assistant/app/config.js}); pirmuosius
šešis perduoda \file{docker-compose.yml}}\label{tab:config:assistant}\\
\toprule
\hd{Kintamasis} & \hd{Numatytoji reikšmė (minimumas)} & \hd{Ką nustato} \\
\midrule
\endfirsthead
\toprule
\hd{Kintamasis} & \hd{Numatytoji reikšmė (minimumas)} & \hd{Ką nustato} \\
\midrule
\endhead
\bottomrule
\endlastfoot
\code{DJANGO_URL} & \code{http://legalops-django:8000} & API: jos viešieji maršrutai, kviečiami su naudotojo prieigos raktu, ir vidiniai maršrutai. \textit{Compose} įrašo tą pačią reikšmę. \\
\code{ASSISTANT_INTERNAL_SECRET} & tuščias & Slaptas. Tas pats bendras slaptas raktas, kaip ir API; kai tuščias, atmetamas kiekvienas vidinis iškvietimas, o paslauga paleidžiama su įspėjimu. \\
\code{AI_GATEWAY_URL} & \code{https://ai.knf.vu.lt/v1} & Fakulteto DI tinklų sąsaja (angl. \en{gateway}; OpenAI duomenų formatai). \textit{Compose} perduoda jį iš \file{.env} su ta pačia numatytąja reikšme. \\
\code{AI_GATEWAY_KEY} & tuščias & Slaptas. DI tinklų sąsajos raktas; kai tuščias, sąsaja nesukonfigūruota: pokalbis atsako 503 (\code{NOT_CONFIGURED}), o generavimo ar vektorizavimo iškvietimas sąsajos modeliui atmetamas su 400 (\code{UNKNOWN_MODEL}). \\
\code{AI_CHAT_MODEL} & \code{gateway:azure/gpt-5.5} & Modelis pokalbio ėjimui ar generavimo iškvietimui, kurio nei instrukcijos versija, nei užklausa modelio nenurodo; visos pateikiamos instrukcijos savąjį nurodo. \textit{Compose} perduoda jį iš \file{.env} su numatytąja reikšme \code{azure/gpt-5.5} (identifikatorius be tiekėjo priešdėlio yra DI tinklų sąsajos identifikatorius). \\
\code{AI_EMBED_MODEL} & \code{gateway:azure/text-embedding-3-large} & Vektorizavimo modelis užklausai, kuri jo nenurodo, --- teisės aktų bazės ar įmonės dokumentų paieškos užklausos vektoriui. Jis turi nurodyti tą patį modelį, kaip ir API kopija (\cref{sec:config:other}). \\
\code{PORT} & \code{3000} (1) & Prievadas, kurio klausosi paslauga. \\
\code{ANTHROPIC_API_KEY} & tuščias & Slaptas, neprivalomas: įgalina modelių identifikatorius \code{anthropic:<model>}. \\
\code{OPENAI_API_KEY} & tuščias & Slaptas, neprivalomas: įgalina modelių identifikatorius \code{openai:<model>}. \\
\code{OLLAMA_BASE_URL} & tuščias & Neprivalomas: Ollama serverio adresas įgalina modelių identifikatorius \code{ollama:<model>}; iš izoliuoto tinklo pasiekiamas tik išėjimo leidžiamų adresų sąrašas. \\
\code{CHAT_TIMEOUT_MS} & \code{900000} (1) & Ilgiausias laikas, kiek gali būti srautu perduodamas visas pokalbio ėjimas. \\
\code{CHAT_STEP_TIMEOUT_MS} & \code{300000} (1) & Ilgiausias laikas, kiek gali trukti vienas ėjimo modelio žingsnis (vienas įrankių ratas). \\
\code{CHAT_CHUNK_TIMEOUT_MS} & \code{90000} (1) & Ilgiausia tyla tarp dviejų srauto dalių; samprotaujantis modelis prieš pirmąjį žetoną galvoja tyliai. \\
\code{GENERATE_TIMEOUT_MS} & \code{900000} (1) & Vieno \route{POST}{/internal/llm/generate} iškvietimo (konvejerio etapo iškvietimo) viršutinė riba --- penkiolika minučių. API laukia 30~s ilgiau (\code{GENERATE_TIMEOUT_S}, \cref{tab:config:calls}). \\
\code{EMBED_TIMEOUT_MS} & \code{60000} (1) & Vieno vektorizavimo paketo viršutinė riba, įskaitant pakartojimus: pusė tų 120~s, kurias laukia API, kad pakibęs iškvietimas taptų klaida, kurią API pakartoja. \\
\code{RATE_LIMIT_TURNS} & \code{12} (1) & Kiek pokalbio ėjimų vienas naudotojas gali pradėti per \code{RATE_LIMIT_WINDOW_MS} (60~s) langą. \\
\code{IDENTITY_CACHE_MS} & \code{15000} (0) & Kiek laiko patikrintu prieigos raktu pasitikima, kol API vėl paklausiama; \code{0} --- klausiama kiekviename ėjime. \\
\code{SHUTDOWN_GRACE_MS} & \code{190000} (0) & Kiek laiko vykstantiems pokalbio ėjimams leidžiama užsibaigti po sustabdymo signalo, kol jie išsaugomi kaip nutraukti. \textit{Compose} nenustato sustabdymo atidėjimo, todėl Docker procesą užbaigia po 10~s (\cref{ch:ops}). \\
\end{xltabular}


\section{Kiti konteineriai}\label{sec:config:other}

Likę konteineriai kintamuosius gauna iš \textit{compose} failo, ir nė vienas jų nesikreipia į
modelį. Jie čia įvardijami išsamumo dėlei, be reikšmių. Įėjimo taškas (Caddy) iš \file{.env}
gauna \code{DBGATE_AUTH_HEADER} --- duomenų bazės įrankio prisijungimo duomenis, užkoduotus HTTP
antraštei. Duomenų bazės konteineriui \textit{compose} faile perduodami \code{POSTGRES_DB},
\code{POSTGRES_USER}, \code{POSTGRES_PASSWORD} ir \code{PGDATA}, o serverio parametrai ---
komandine eilute (tarp jų \code{max_connections=100}). Duomenų bazės įrankiui (dbgate)
perduodamas jo žiniatinklio šakninis kelias, pagrindinio autentifikavimo prisijungimas su
\code{PASSWORD} iš \file{.env} (\code{DBGATE_PASSWORD}) ir vienas ryšys su duomenų baze.
Konsolė, išėjimo tarpinis serveris ir užduočių vykdytojo konteineris (kuris tik paleidžia
užduočių vykdytoją API konteineryje) aplinkos negauna. Kaip slapti duomenys sukuriami ir
keičiami, aprašyta \cref{ch:ops}.

Vienas kintamasis nustatomas du kartus, ir abi kopijos turi sutapti. API \code{AI_EMBED_MODEL}
įrašytas \textit{compose} faile ir nurodo modelį kiekviename iškvietime, kuris išsaugo vektorių;
asistento paslaugos kopija ateina iš \file{.env} ir yra kiekvieno užklausos vektoriaus modelis.
Dviejų modelių vektorių palyginti negalima, todėl, jei jie skirtųsi, paieška rikiuotų išsaugotus
vektorius pagal kito modelio užklausos vektorių. Kiekvieno vektoriaus ilgis fiksuotas --- 1536
(\code{EMBED_DIMENSIONS}).


\section{Derinamos konstantos}\label{sec:config:constants}

Šio skyriaus lentelėse išvardytos konstantos, kurios lemia, kiek modelis perskaito per vieną
iškvietimą, kiek iškvietimų vykdoma vienu metu, kiek laiko iškvietimas gali trukti ir kiek kartų
bandomas, ką kodas priima iš atsakymo ir kokiu raktu rezultatas saugomas podėlyje. Reikšmės
pateikiamos taip, kaip jas rašo kodas (\code{24_000} reiškia dvidešimt keturis tūkstančius). Kai
riba įrašyta tiesiai funkcijoje, o ne įvardyta, lentelėje ji žymima \enquote{(įterpta)}. Failų
keliai yra API kodo, nebent baigiasi \file{.js} --- tada tai asistento paslaugos failai. Konstantos,
kurios tik išvardija leistinas reikšmes (būsenos, reikšmingumai, kategorijos, kurias apibrėžia
schemos išvardijimas), neišvardijamos: jos yra kiekvieno DI kelio schemos dalis ir aprašytos kartu su
tuo keliu.


\subsection{Modelių iškvietimai}\label{sec:config:calls}

Šios konstantos lemia kiekvieną iškvietimą, kad ir kokia būtų jo paskirtis: kiek jų vykdoma
kartu, kaip jie kartojami, kiek laukia kiekviena pusė ir ką gali talpinti vektorizavimo paketas.

\begin{lessonbox}{Išmatuota 2026~m. lapkričio 10~d.}
Keturi vienu metu vykdyti atitikties spragų analizės iškvietimai, kiekvienas maždaug po
59\,000 žetonų, per dvi minutes pasiekė tiekėjo žetonų per minutę ribą, o įprasti pakartojimai
po 2, 6 ir 18~sekundžių ją pasiekė vėl. Kodo komentaras tai nurodo kaip priežastį, kodėl
užduoties etapui priklausantis iškvietimas yra \enquote{kantrus}: užklausų dažnio ribos
iš viso laukiama iki penkiolikos minučių (\code{RATE_LIMIT_PATIENCE_S}), užuot leidus etapui
nepavykti. Visa 2026~m. lapkričio 13~d. ataskaita (272 iškvietimai) dėl ribos sulėtėjo, bet
neprarado nė vieno iškvietimo.
\end{lessonbox}

\begin{xltabular}{\textwidth}{@{}L{44mm}L{19mm}Y@{}}
\caption{Modelių iškvietimų konstantos}\label{tab:config:calls}\\
\toprule
\hd{Konstanta, failas} & \hd{Reikšmė} & \hd{Poveikis} \\
\midrule
\endfirsthead
\toprule
\hd{Konstanta, failas} & \hd{Reikšmė} & \hd{Poveikis} \\
\midrule
\endhead
\bottomrule
\endlastfoot
\code{WORKERS}\newline{\footnotesize\color{muted}\file{analysis/pipeline.py}} & \code{4} & Kiek ataskaitos etapo modelio iškvietimų vykdoma vienu metu. \\
\code{shared_first}\newline{\footnotesize\color{muted}\file{analysis/pipeline.py}} & -- & Ne skaičius, o taisyklė: iš etapo iškvietimų, kurie dalijasi duomenimis, pirmasis vykdomas vienas, kad tiekėjas bendrąjį bloką turėtų podėlyje, kol kreipsis kiti. \\
\code{workers=4}\newline{\footnotesize\color{muted}\file{corpus/obligations.py}, \file{corpus/verify.py}, \file{corpus/outline.py} (įterpta)} & \code{4} & Kiek iškvietimų vykdoma vienu metu, kai išskiriamos pareigos, kai vertintojas jas tikrina ir kai skaitomi ilgi standarto punktai; \code{WORKERS = 4} faile \file{programme/run.py} tą patį daro programos dokumento inventoriams. \\
\code{patient}\newline{\footnotesize\color{muted}\file{assistant/llm.py} (įterpta)} & -- & \code{run_prompt} iškvietimas, atliekamas užduoties etape (nurodytas \code{job_step_id}), užklausų dažnio ribos palaukia; iškvietimas, kurio žmogus laukia prie ekrano, nelaukia. \\
\code{BACKOFF_BASE_S}, \code{BACKOFF_MAX_S}\newline{\footnotesize\color{muted}\file{assistant/gateway.py}} & \code{2}, \code{30} & Sekundės iki $n$-ojo pakartojimo: $\min(2\cdot 3^n, 30)$, t.~y. 2, 6, 18, 30. \\
\code{retries=1}\newline{\footnotesize\color{muted}\file{assistant/gateway.py} (\code{generate} argumentas)} & \code{1} & Modelio iškvietimas, nepavykęs su kartotina būsena, bandomas dar kartą. \\
\code{CONNECTION_RETRIES}\newline{\footnotesize\color{muted}\file{assistant/gateway.py}} & \code{5} & Generavimo ar vektorizavimo iškvietimas, kuris išvis nepasiekė asistento paslaugos (paslauga persikrauna), kartojamas penkis kartus, po 2, 6, 18, 30 ir 30~s; dėl laiko ribos nutrūkęs iškvietimas išlaiko įprastus pakartojimus. \\
\code{RETRYABLE_STATUSES}\newline{\footnotesize\color{muted}\file{assistant/gateway.py}} & 408, 409, 425, 429, 500, 502, 503, 504 & HTTP būsenos, kurių atveju, kaip ir tinklo gedimų atveju, iškvietimas kartojamas; dėl bet kurios kitos klaidos iškvietimas iškart nepavyksta. \\
\code{RATE_LIMIT_WAIT_S}\newline{\footnotesize\color{muted}\file{assistant/gateway.py}} & \code{20} & Kantraus iškvietimo laukimas, kai tiekėjas jo nenurodo: 20~s, padaugintos iš laukimo eilės numerio, ne daugiau kaip 60~s, atsitiktinai ketvirtadaliu daugiau ar mažiau, kad etapo iškvietimai negrįžtų kartu. \\
\code{RATE_LIMIT_PATIENCE_S}\newline{\footnotesize\color{muted}\file{assistant/gateway.py}} & \code{900} & Kiek iš viso sekundžių kantrus iškvietimas laukia, kol praeis užklausų dažnio ribos, prieš nepavykdamas. \\
\code{GENERATE_TIMEOUT_S}\newline{\footnotesize\color{muted}\file{assistant/gateway.py}} & \code{930} & Kiek API laukia vieno generavimo iškvietimo: 30~s ilgiau už asistento paslaugos \code{GENERATE_TIMEOUT_MS}, kad pirmiau ateitų pačios paslaugos atsisakymas. \\
\code{timeout=120}, \code{retries=3}\newline{\footnotesize\color{muted}\file{assistant/gateway.py} (\code{embed} argumentai)} & \code{120}, \code{3} & Kiek sekundžių API laukia vieno vektorizavimo paketo, ir kiek kartų jis kartojamas. \\
\code{MAX_EMBED_TEXTS}\newline{\footnotesize\color{muted}\file{routes/internal.js}} & \code{256} & Daugiausia tekstų, kuriuos gali nešti vienas vektorizavimo iškvietimas. \\
\code{EMBED_BATCH}\newline{\footnotesize\color{muted}\file{corpus/ingest.py}} & \code{64} & Teisės akto fragmentų skaičius viename vektorizavimo iškvietime. Pareigos siunčiamos po 32 eilutes per iškvietimą, kiekviena su lietuvišku ir anglišku tekstu (\file{corpus/obligations.py}, įterpta), o įmonės dokumento ištraukos --- po 64 (\file{companies/indexing.py}, įterpta). \\
\code{EMBED_DIMENSIONS}\newline{\footnotesize\color{muted}\file{corpus/models.py}, \file{companies/models.py}, \file{routes/internal.js}} & \code{1536} & Kiekvieno išsaugoto vektoriaus ilgis; \code{text-embedding-3} modelio prašoma būtent jo, o ilgesnis vektorius nukerpamas iki pirmųjų 1536 reikšmių ir vėl normalizuojamas. \\
\code{MODELS_CACHE_MS}\newline{\footnotesize\color{muted}\file{routes/internal.js}} & \code{10 * 60_000} & Kiek laiko pakartotinai naudojamas DI tinklų sąsajos modelių sąrašas (su kainomis). \\
\code{PRICE_TTL_S}, \code{PRICE_RETRY_S}\newline{\footnotesize\color{muted}\file{assistant/usage.py}} & \code{3600}, \code{300} & Kainų lentelė, pagal kurią įkainojamas kiekvienas iškvietimas, atnaujinama kas valandą; nepavykus atnaujinti, naudojamos išsaugotos kainos, o atnaujinti bandoma vėl po penkių minučių. \\
\code{MIN_BUDGET_TO_START_USD}\newline{\footnotesize\color{muted}\file{analysis/api/views.py}, \file{programme/api/views.py}} & \code{0.25} & Kiek dolerių paskyrai šiandien turi būti likę, kad galėtų pradėti ataskaitą, patvirtinti ataskaitos apimtį ar pradėti programos dokumentą; jau pradėtas vykdymas dėl biudžeto niekada nestabdomas. \\
\end{xltabular}


\subsection{Teisės aktų bazė}\label{sec:config:corpus}

Šios konstantos lemia, kaip teisės aktas ar standartas suskaidomas į vienetus ir fragmentus, kaip
ieškoma teisės aktų bazėje, kaip modelis perskaito standarto, kurio numeracijos taisyklės negali
sekti, struktūrą ir kaip nuorodų patikra perskaito nurodytą nuostatą. Kiekvieną mechanizmą
paaiškina \cref{ch:corpus}.

\begin{xltabular}{\textwidth}{@{}L{44mm}L{19mm}Y@{}}
\caption{Teisės aktų bazės konstantos}\label{tab:config:corpus}\\
\toprule
\hd{Konstanta, failas} & \hd{Reikšmė} & \hd{Poveikis} \\
\midrule
\endfirsthead
\toprule
\hd{Konstanta, failas} & \hd{Reikšmė} & \hd{Poveikis} \\
\midrule
\endhead
\bottomrule
\endlastfoot
\code{MAX_CHUNK_CHARS}\newline{\footnotesize\color{muted}\file{corpus/chunking.py}} & \code{3200} & Ilgiausias paieškos fragmentas (apie 800 žetonų); ilgesnis straipsnis skaidomas ties dalių ribomis. \\
\code{SHORT_CLAUSE_CHARS}, \code{RUN_CHARS}\newline{\footnotesize\color{muted}\file{corpus/chunking.py}} & \code{700}, \code{2000} & Trumpesnis nei 700 simbolių standarto punktas dalijasi fragmentu su trumpais gretimais punktais, iki 2000. \\
\code{RRF_K}\newline{\footnotesize\color{muted}\file{corpus/search.py}} & \code{60} & Abipusio rango suliejimo (RRF) konstanta, kuria sujungiamos vektorinė ir raktažodžių paieškos pusės. \\
\code{CANDIDATES}\newline{\footnotesize\color{muted}\file{corpus/search.py}} & \code{40} & Kiek fragmentų kiekviena paieškos pusė pateikia prieš sujungimą. \\
\code{WINDOW_CHARS}\newline{\footnotesize\color{muted}\file{corpus/outline.py}} & \code{28_000} & Kiek standarto teksto modelis perskaito per vieną iškvietimą, kai žymi antraštes (apie 8\,000 žetonų), sunumeruotomis eilutėmis; vieno dokumento ištraukos skaitomos viena po kitos. \\
\code{SHOWN_LINE_CHARS}\newline{\footnotesize\color{muted}\file{corpus/outline.py}} & \code{200} & Kiekviena eilutė modeliui rodoma iki šio ilgio. \\
\code{CONTEXT_HEADINGS}\newline{\footnotesize\color{muted}\file{corpus/outline.py}} & \code{14} & Kiek prieš ištrauką rastų antraščių pasakoma modeliui, kad jis išlaikytų lygius. \\
\code{MAX_AI_CHARS}\newline{\footnotesize\color{muted}\file{corpus/outline.py}} & \code{2_500_000} & Ilgesnis dokumentas modeliui niekada nesiunčiamas; jis perskaitomas taisyklėmis arba pagal puslapius. \\
\code{REFINE_FROM_CHARS}, \code{MIN_REFINED}\newline{\footnotesize\color{muted}\file{corpus/outline.py}} & \code{15_000}, \code{3} & Punktą, kurį taisyklės paliko ilgesnį nei 15\,000 simbolių, modelis skaito toliau, o jo pažymėtos antraštės išlaikomos, kai bent trys turi identifikatorių. \\
\code{MAX_LEAF_SHARE}\newline{\footnotesize\color{muted}\file{corpus/outline.py}} & \code{0.5} & Kai vienas taisyklių paliktas ilgesnis nei 15\,000 simbolių punktas apima daugiau nei pusę teksto, modelio prašoma perskaityti visą dokumentą. \\
\code{SPLIT_FROM_CHARS}, \code{PART_CHARS}\newline{\footnotesize\color{muted}\file{corpus/outline.py}} & \code{24_000}, \code{6_000} & Vienetas, vis dar ilgesnis nei 24\,000 simbolių, karpomas pagal puslapius, o kai tekste puslapių nesimato --- į dalis po 6\,000, kad nė vienas reikalavimas nebūtų cituojamas nurodant keturiasdešimt puslapių. \\
\code{MIN_STANDARD_CLAUSES}, \code{MAX_PREAMBLE_SHARE}\newline{\footnotesize\color{muted}\file{corpus/standards.py}} & \code{3}, \code{0.6} & Numeracijos taisyklės perskaito standartą, kai randa bent tris punktus, o prieš pirmąjį yra ne daugiau kaip 60~\% teksto; kitaip modelio prašoma pažymėti antraštes, o nepavykus vietoj jų naudojami puslapiai. \\
\code{OWN_TEXT_CHARS}\newline{\footnotesize\color{muted}\file{corpus/standards.py}} & \code{40} & Paties punkto tekstas prieš jo papunkčius skaitomas ieškant reikalavimų, kai yra bent tokio ilgio. \\
\code{MAX_ITEMS}\newline{\footnotesize\color{muted}\file{corpus/citations.py}} & \code{50} & Kiek nuorodų įvertina viena nuorodų patikra (vienas modelio iškvietimas); likusios ilgesnės užklausos nuorodos netikrinamos. \\
\code{MAX_LOOKUPS}\newline{\footnotesize\color{muted}\file{corpus/citations.py}} & \code{2 * MAX_ITEMS} & Kiek straipsnių paieškų gali atlikti viena nuorodų patikros užklausa. \\
\code{EXCERPT_CHARS}\newline{\footnotesize\color{muted}\file{corpus/citations.py}} & \code{600} & Nuostatos pradžia, su kuria vertintojas lygina teiginį (kartu su jos antrašte). \\
\code{MAX_SERIES_ARTICLES}\newline{\footnotesize\color{muted}\file{corpus/citations.py}} & \code{15} & Kiek vieno intervalo ar sąrašo (\enquote{6--9, 26--29 str.}) straipsnių ieškoma po vieną. \\
teiginys, pastaba (įterpta)\newline{\footnotesize\color{muted}\file{corpus/citations.py}} & \code{500}, \code{300} & Teiginys nukerpamas iki 500 simbolių, prieš jį perskaitant vertintojui; vertintojo pastaba apribojama 300 simbolių. \\
\end{xltabular}

Faile \file{corpus/standards.py} yra dar maždaug dvidešimt skaitinių ribų ir tiek pat teksto
šablonų --- kiek eilučių puslapio krašte gali būti puslapio apipavidalinimas, kokio ilgio gali
būti antraštė, kaip atpažįstama perkelta antraštė ar dalies numeris. Jos lemia tik tai, kaip
taisyklės perskaito standarto tekstą, niekada --- ką gauna modelis, ir paaiškintos ten, kur
įrašytos tame faile.


\subsection{Pareigos, apimtis ir sankcijos}\label{sec:config:obligations}

Trys teisės aktų bazės skaitymai, kuriuos modelis atlieka vieną kartą kiekvienai teisės akto
redakcijai, ir vertintojas, kuris tikrina registrą (\cref{ch:obligations}).

\begin{xltabular}{\textwidth}{@{}L{44mm}L{19mm}Y@{}}
\caption{Pareigų, apimties skaitymo, sankcijų ir vertintojo
konstantos}\label{tab:config:obligations}\\
\toprule
\hd{Konstanta, failas} & \hd{Reikšmė} & \hd{Poveikis} \\
\midrule
\endfirsthead
\toprule
\hd{Konstanta, failas} & \hd{Reikšmė} & \hd{Poveikis} \\
\midrule
\endhead
\bottomrule
\endlastfoot
\code{BATCH_CHARS}\newline{\footnotesize\color{muted}\file{corpus/obligations.py}} & \code{24_000} & Nuostatų simbolių skaičius viename išskyrimo iškvietime (apie 6\,000 žetonų). \\
\code{MAX_PROVISION_CHARS}\newline{\footnotesize\color{muted}\file{corpus/obligations.py}} & \code{60_000} & Ilgesnė pavienė nuostata savo pakete nukerpama. \\
\code{WHOLE_PART_CHARS}\newline{\footnotesize\color{muted}\file{corpus/obligations.py}} & \code{400} & Skyrius ar skirsnis, kurio tekstas visai nenumeruotas, perskaitomas visas, kai yra ilgesnis už šią ribą. \\
\code{DUPLICATE_RATIO}\newline{\footnotesize\color{muted}\file{corpus/obligations.py}} & \code{0.75} & Pavadinimų panašumas, nuo kurio naujas pasiūlymas, cituojantis tą pačią nuostatą, laikomas patvirtintos pareigos dublikatu ir atmetamas. \\
\code{BATCH_CHARS}\newline{\footnotesize\color{muted}\file{corpus/verify.py}} & \code{24_000} & Straipsnių simbolių skaičius, kurį vertintojas perskaito per vieną iškvietimą. \\
\code{PROVISION_CHARS}\newline{\footnotesize\color{muted}\file{corpus/verify.py}} & \code{20_000} & Kiekvienas vertintojui rodomas straipsnis nukerpamas iki šio ilgio ir pažymimas \enquote{nukirpta}. \\
\code{OBLIGATIONS_PER_CALL}\newline{\footnotesize\color{muted}\file{corpus/verify.py}} & \code{20} & Kiek pareigų vertintojas patikrina per vieną iškvietimą. \\
\code{ACTS_PER_CALL}\newline{\footnotesize\color{muted}\file{corpus/scope.py}} & \code{4} & Kiek teisės aktų taikymo srities nuostatas perskaito vienas iškvietimas. \\
\code{SCOPE_ARTICLES}\newline{\footnotesize\color{muted}\file{corpus/scope.py}} & \code{3} & Pirmieji teisės akto straipsniai, kurie lemia jo apimtį; aktas be straipsnių pateikia pirmuosius šešis punktus (įterpta). \\
\code{READ_CHARS}\newline{\footnotesize\color{muted}\file{corpus/scope.py}} & \code{12_000} & Kiekvienos taikymo srities nuostatos simbolių skaičius apimties skaitymo metu. \\
\code{HEADING}\newline{\footnotesize\color{muted}\file{corpus/sanctions.py}} & šablonas & Straipsnis laikomas sankcijų nuostata, kai jo antraštėje lietuviškai ar angliškai minimos sankcijos, baudos, nuobaudos ar priemonės. \\
\code{ARTICLES}\newline{\footnotesize\color{muted}\file{corpus/sanctions.py}} & \code{8} & Kiek vieno teisės akto sankcijų nuostatų perskaitoma per vienintelį jo iškvietimą. \\
\code{READ_CHARS}\newline{\footnotesize\color{muted}\file{corpus/sanctions.py}} & \code{14_000} & Kiekvienos perskaitomos sankcijų nuostatos simbolių skaičius. \\
\code{MAX_TIERS}\newline{\footnotesize\color{muted}\file{corpus/sanctions.py}} & \code{24} & Kiek pakopų išlaikoma iš vieno atsakymo; instrukcija prašo ne daugiau kaip dvidešimt keturių. \\
\end{xltabular}


\subsection{Dokumentai ir inventorius}\label{sec:config:documents}

Kas perskaitoma iš įmonės dokumentų ir kokiais raktais tai saugoma podėlyje, kad vėlesnė
ataskaita galėtų tai panaudoti dar kartą (\cref{ch:company}).

\begin{xltabular}{\textwidth}{@{}L{44mm}L{19mm}Y@{}}
\caption{Įmonės dokumentų konstantos}\label{tab:config:documents}\\
\toprule
\hd{Konstanta, failas} & \hd{Reikšmė} & \hd{Poveikis} \\
\midrule
\endfirsthead
\toprule
\hd{Konstanta, failas} & \hd{Reikšmė} & \hd{Poveikis} \\
\midrule
\endhead
\bottomrule
\endlastfoot
\code{MAX_DOCUMENT_CHARS}\newline{\footnotesize\color{muted}\file{companies/inventory.py}} & \code{160_000} & Kiek dokumento simbolių perskaito inventoriaus iškvietimas (apie 40\,000 žetonų); modeliui pasakoma, kiek dar liko neperskaityta. \\
\code{SCHEMA_VERSION}\newline{\footnotesize\color{muted}\file{companies/inventory.py}} & \code{2} & Inventoriaus forma, jo podėlio rakto dalis; ją būtina padidinti kiekvieną kartą, kai pasikeičia inventoriaus schema, kad kiekvienas inventorius būtų perskaitytas iš naujo. \\
inventoriaus raktas\newline{\footnotesize\color{muted}\file{companies/inventory.py} (\code{inventory_key})} & -- & \enquote{teksto sha256\,:\,instrukcijos versijos id\,:\,s\textit{SCHEMA\_VERSION}}: inventorius naudojamas pakartotinai, kol nepasikeičia tekstas, instrukcijos versija ar forma. \\
\code{CITATION_FIELDS}\newline{\footnotesize\color{muted}\file{companies/inventory.py}} & aktas 300, straipsnis 60, dalis 60, punktas 60, teiginys 1000, citata 200 & Ilgiausia išlaikoma kiekvieno inventoriaus rastos nuorodos lauko reikšmė; instrukcija prašo ne ilgesnių kaip 200 simbolių citatų. \\
\code{WINDOW_CHARS}, \code{OVERLAP_CHARS}\newline{\footnotesize\color{muted}\file{companies/indexing.py}} & \code{1800}, \code{200} & Ištraukos, į kurias dokumentas suskaidomas vektoriams (apie 450 žetonų); kiekviena iš dalies sutampa su ankstesne. \\
\code{RESUME_AFTER_RESET}\newline{\footnotesize\color{muted}\file{companies/indexing.py}} & 5 minutės & Indeksavimo užduotis, atidėta todėl, kad įkėlusio naudotojo dienos biudžetas išnaudotas, grįžta praėjus tiek laiko po 00:00 UTC. \\
\end{xltabular}


\subsection{Ataskaita}\label{sec:config:report}

Ataskaitos konvejerio (\cref{ch:report}), jo galimybės įvertinti įmonės vidaus reikalavimus ir
prieš ataskaitos pradžią rodomo kainos įverčio konstantos. Įverčio apimtys vienam iškvietimui
išmatuotos su pavyzdine įmone \demo{} (failas pirmąsias tris datuoja 2026~m. lapkričio 10~d., su
\model{gpt-6.1-sol}); tai viršutinės ribos įvestys, o ne prognozės.

\begin{xltabular}{\textwidth}{@{}L{44mm}L{19mm}Y@{}}
\caption{Atitikties ataskaitos konstantos}\label{tab:config:report}\\
\toprule
\hd{Konstanta, failas} & \hd{Reikšmė} & \hd{Poveikis} \\
\midrule
\endfirsthead
\toprule
\hd{Konstanta, failas} & \hd{Reikšmė} & \hd{Poveikis} \\
\midrule
\endhead
\bottomrule
\endlastfoot
\code{CITATIONS_PER_CALL}\newline{\footnotesize\color{muted}\file{analysis/pipeline.py}} & \code{MAX_ITEMS} (50) & Teisės aktų nuorodų patikros paieškos viename vertintojo iškvietime. \\
\code{ACTS_PER_CALL}\newline{\footnotesize\color{muted}\file{analysis/pipeline.py}} & \code{4} & Teisės aktai viename taikytinumo iškvietime. \\
\code{SCOPE_CHARS}\newline{\footnotesize\color{muted}\file{analysis/pipeline.py}} & \code{5000} & Kiekvienos taikymo srities nuostatos simbolių skaičius taikytinumo iškvietime. \\
\code{OBLIGATIONS_PER_CALL}\newline{\footnotesize\color{muted}\file{analysis/pipeline.py}} & \code{15} & Pareigos, vertinamos viename atitikties spragų analizės iškvietime; visos iš vieno teisės akto. \\
\code{DOCUMENT_PACK_CHARS}\newline{\footnotesize\color{muted}\file{analysis/pipeline.py}} & \code{600_000} & Kiek įmonės dokumentų simbolių gali nešti vienas iškvietimas. Į šią ribą telpantis dokumentų rinkinys siunčiamas visas, kaip bendri duomenys, kiekvienam atitikties spragų analizės iškvietimui (ir vidaus reikalavimų iškvietimams bei programos dokumentams); didesnis apkarpomas iki dokumentų, artimiausių kiekvienam paketui, o vidaus reikalavimų patikra ir programos dokumentai tada nerengiami. \\
\code{EVIDENCE_PER_OBLIGATION}, \code{EVIDENCE_MAX_DISTANCE}\newline{\footnotesize\color{muted}\file{analysis/pipeline.py}} & \code{6}, \code{0.7} & Per dideliam dokumentų rinkiniui: šešios pareigai artimiausios ištraukos, kurių kosinusinis atstumas mažesnis nei 0,7, pagal geresnį iš lietuviško ir angliško vektoriaus, lemia, kuriuos dokumentus gauna paketas. \\
\code{PLAN_FINDINGS_PER_CALL}\newline{\footnotesize\color{muted}\file{analysis/pipeline.py}} & \code{150} & Kiek atvirų išvadų paima vienas veiksmų plano iškvietimas; ištisos temos sudedamos kartu. \\
\code{MAX_CITATIONS_FOR_PLAN}\newline{\footnotesize\color{muted}\file{analysis/pipeline.py}} & \code{60} & Daugiausia teisės aktų nuorodų patikros problemų, kurias perskaito veiksmų planas. \\
plano eilutės (įterpta)\newline{\footnotesize\color{muted}\file{analysis/pipeline.py}} & \code{280}, \code{220}, \code{600} & Išvada planavimo iškvietimui pateikiama su analize, nukirpta iki 280 simbolių, ir rekomendacija --- iki 220; veiksmo projektas sujungimo iškvietimui --- su aprašymu, nukirptu iki 600. \\
santraukos įvestis (įterpta)\newline{\footnotesize\color{muted}\file{analysis/pipeline.py} (\code{step_summary})} & \code{25}, \code{5}, \code{30} & Santraukai pateikiama ne daugiau kaip 25 rimčiausios atviros išvados (neatitinka arba iš dalies, kritinės arba aukštos), papildytos iki penkių mažiau reikšmingomis, penkios didžiausios vidaus reikalavimų spragos ir pirmieji 30 plano veiksmų. \\
\code{REQUIREMENTS_PER_DOCUMENT}\newline{\footnotesize\color{muted}\file{analysis/internal.py}} & \code{20} & Kiek vidaus reikalavimų išlaikoma iš vieno politikos dokumento; instrukcija prašo ne daugiau kaip dvidešimties. \\
\code{READ_CHARS}\newline{\footnotesize\color{muted}\file{analysis/internal.py}} & \code{80_000} & Kiek politikos dokumento simbolių perskaitoma ieškant jo reikalavimų. \\
\code{PER_CALL}\newline{\footnotesize\color{muted}\file{analysis/internal.py}} & \code{15} & Kiek vidaus reikalavimų įvertinama per vieną iškvietimą. \\
reikalavimų raktas\newline{\footnotesize\color{muted}\file{analysis/internal.py} (įterpta)} & -- & \enquote{įkelto failo sha256\,:\,instrukcijos versijos id}: dokumentas reikalavimams perskaitomas iš naujo tik tada, kai pasikeičia vienas iš jų. \\
\code{APPLICABILITY_CALL_TOKENS}\newline{\footnotesize\color{muted}\file{analysis/estimate.py}} & \code{(24_000, 1_800)} & Keturių teisės aktų taikytinumo iškvietimo įvesties ir išvesties žetonai. \\
\code{INVENTORY_CALL_TOKENS}\newline{\footnotesize\color{muted}\file{analysis/estimate.py}} & \code{(2_800, 7_000)} & Tas pats vieno dokumento inventoriui. \\
\code{CITATION_CALL_TOKENS}\newline{\footnotesize\color{muted}\file{analysis/estimate.py}} & \code{(13_000, 3_500)} & Tas pats vienai nuorodų patikrai. \\
\code{GAP_CALL_TOKENS}, \code{GAP_PREFIX_TOKENS}\newline{\footnotesize\color{muted}\file{analysis/estimate.py}} & \code{(5_500, 4_500)}, \code{5_000} & Ką prideda ir ką parašo penkiolikos pareigų atitikties spragų analizės iškvietimas, ir profilis bei instrukcija prieš dokumentus. \\
\code{DOCUMENT_CHARS_PER_TOKEN}\newline{\footnotesize\color{muted}\file{analysis/estimate.py}} & \code{4.5} & Įmonės dokumento simbolių skaičius vienam žetonui. \\
\code{PLAN_CALL_TOKENS}, \code{SUMMARY_TOKENS}\newline{\footnotesize\color{muted}\file{analysis/estimate.py}} & \code{(31_000, 5_000)}, \code{(14_000, 2_000)} & Vienas planavimo iškvietimas ir santrauka. \\
\code{OPEN_SHARE}\newline{\footnotesize\color{muted}\file{analysis/estimate.py}} & \code{0.5} & Kokia vertinamų pareigų dalis, kaip manoma, baigsis atviromis išvadomis ir pateks į planą. \\
\code{CHARS_PER_TOKEN}, \code{EXTRACTION_OUTPUT_RATIO}\newline{\footnotesize\color{muted}\file{analysis/estimate.py}} & \code{3.2}, \code{0.7} & Teisės aktas, iš kurio pareigos dar niekada neišskirtos: įvesties žetonai --- jo simbolių skaičius, padalytas iš 3,2, išvestis --- 0,7 įvesties. \\
\code{OBLIGATIONS_PER_ARTICLE}\newline{\footnotesize\color{muted}\file{analysis/estimate.py}} & \code{2.2} & Kiek pareigų numanoma vienam straipsniui teisės akto, iš kurio pareigos dar neišskirtos. \\
\code{CITATIONS_PER_DOCUMENT}\newline{\footnotesize\color{muted}\file{analysis/estimate.py}} & \code{20} & Kiek nuorodų numanoma dokumente, kurio inventorius bus perskaitytas iš naujo. \\
\code{EXPECTED_PER_DOCUMENT}\newline{\footnotesize\color{muted}\file{analysis/internal.py}} & \code{8} & Kiek vidaus reikalavimų numanoma dar neperskaitytame politikos dokumente. \\
\end{xltabular}


\subsection{Veiksmų registras}\label{sec:config:register}

Registras (\cref{ch:register}) modelio nekviečia; dvi jo konstantos lemia, ką mato žmogus.

\begin{xltabular}{\textwidth}{@{}L{44mm}L{19mm}Y@{}}
\caption{Veiksmų registro konstantos}\label{tab:config:register}\\
\toprule
\hd{Konstanta, failas} & \hd{Reikšmė} & \hd{Poveikis} \\
\midrule
\endfirsthead
\toprule
\hd{Konstanta, failas} & \hd{Reikšmė} & \hd{Poveikis} \\
\midrule
\endhead
\bottomrule
\endlastfoot
\code{DUE_DAYS}\newline{\footnotesize\color{muted}\file{analysis/register.py}} & 14, 30, 90, 365 & Dienos nuo patvirtinimo iki veiksmo termino laikotarpiams \code{immediate}, \code{30_days}, \code{90_days} ir \code{12_months}; bet kuriam kitam --- 90. \\
\code{HISTORY}\newline{\footnotesize\color{muted}\file{analysis/register.py}} & \code{4} & Kiek ankstesnių vertinimų apima išvados istorija. \\
\code{OPEN}, \code{SETTLED}, \code{LIVE}\newline{\footnotesize\color{muted}\file{analysis/register.py}} & -- & Kurios išvadų būsenos laikomos atviromis (neatitinka, iš dalies, neįrodyta) ar išspręstomis (atitinka, netaikoma) ir kurios veiksmų būsenos yra aktyvios (\code{open}, \code{in_progress}, \code{done}). \\
\end{xltabular}


\subsection{Programos dokumentai}\label{sec:config:programme}

Programos dokumentas (\cref{ch:programme}) įmonę perskaito taip pat, kaip ataskaita, --- ta pati
\code{DOCUMENT_PACK_CHARS}, ta pati nuorodų patikra, --- o paskui atsakymą tikrina jo
taisyklės. Blokuojanti taisyklė neleidžia dokumento patvirtinti.

\begin{xltabular}{\textwidth}{@{}L{44mm}L{19mm}Y@{}}
\caption{Programos dokumentų konstantos}\label{tab:config:programme}\\
\toprule
\hd{Konstanta, failas} & \hd{Reikšmė} & \hd{Poveikis} \\
\midrule
\endfirsthead
\toprule
\hd{Konstanta, failas} & \hd{Reikšmė} & \hd{Poveikis} \\
\midrule
\endhead
\bottomrule
\endlastfoot
\code{STARTS}\newline{\footnotesize\color{muted}\file{programme/run.py}} & \code{(0, 5, 50, 92)} & Kur procentais prasideda kiekvienas iš keturių vykdymo etapų --- eigai, rodomai, kol žmogus laukia. \\
\code{PART_CHARS}\newline{\footnotesize\color{muted}\file{programme/policy.py}} & \code{12_000} & Ilgiausias išlaikomas vienos politikos dalies tekstas. \\
\code{MIN_PART_CHARS}\newline{\footnotesize\color{muted}\file{programme/policy.py}} & \code{150} & Trumpesnė dalis pažeidžia blokuojančią taisyklę, kad visos aštuonios dalys turi būti parašytos. \\
\code{MIN_STEPS}\newline{\footnotesize\color{muted}\file{programme/governance.py}} & \code{5} & Mažiausias atitikties proceso žingsnių skaičius (blokuojanti taisyklė); instrukcija prašo nuo aštuonių iki keturiolikos. \\
\code{ISO_CLAUSES}\newline{\footnotesize\color{muted}\file{programme/governance.py}} & 41 punktas & ISO 37301:2021 punktai, kuriais gali remtis proceso žingsnis. \\
\code{GUIDELINE_PARAGRAPHS}\newline{\footnotesize\color{muted}\file{programme/governance.py}} & \code{range(1, 80)} & Lietuvos banko gairių punktai (nuo 1 iki 79), kuriais gali remtis žingsnis. \\
\code{MIN_SCENARIOS}, \code{MIN_CATEGORIES}\newline{\footnotesize\color{muted}\file{programme/risks.py}} & \code{5}, \code{4} & Mažiausias scenarijų ir teisinės rizikos rūšių skaičius registre (abi taisyklės blokuojančios); instrukcija prašo nuo šešių iki dešimties, apimančių bent penkias rūšis. \\
\code{BANDS}\newline{\footnotesize\color{muted}\file{programme/risks.py}} & 1--3, 4--6, 8--12, 15--16, 20--25 & Penkios 5×5 matricos zonos ir kiekvienoje leidžiami atsakai; laikomos kode, modelio niekada neklausiama. \\
\code{DECISION_CHARS}, \code{CALIBRATION_CHARS}\newline{\footnotesize\color{muted}\file{programme/risks.py}} & \code{1_000}, \code{4_000} & Ilgiausias išlaikomas zonos sprendimas ir ilgiausia išlaikoma pastaba apie tai, kaip pritaikytos skalės: vietos palikta tam, ko prašo instrukcija, kad nė vienas tekstas nebūtų nukirptas žodžio viduryje. \\
\end{xltabular}


\subsection{Pokalbis}\label{sec:config:chat}

Pokalbio ėjimo ribos, ribos to, ką modelis gauna apie naudotojo įmonę, ir išorinių įrankių,
kuriuos gali prijungti administratorius, ribos (\cref{ch:chat}). Kiekvieno pokalbio įrankio
grąžinamų duomenų ribos įrašytos įrankių failuose kataloge \file{llm/tools/} ir išvardytos kartu su
įrankiais (\cref{app:tools}). Pokalbio aplinkos kintamieji pateikti \cref{tab:config:assistant}.

\begin{xltabular}{\textwidth}{@{}L{44mm}L{19mm}Y@{}}
\caption{Pokalbių asistento konstantos (asistento paslauga)}\label{tab:config:chat}\\
\toprule
\hd{Konstanta, failas} & \hd{Reikšmė} & \hd{Poveikis} \\
\midrule
\endfirsthead
\toprule
\hd{Konstanta, failas} & \hd{Reikšmė} & \hd{Poveikis} \\
\midrule
\endhead
\bottomrule
\endlastfoot
\code{MAX_STEPS}\newline{\footnotesize\color{muted}\file{config.js}} & \code{24} & Kiek modelio žingsnių gali sujungti vienas ėjimas; paskutinis vykdomas be įrankių, kad ėjimas baigtųsi tekstu. \\
\code{MAX_MESSAGE_CHARS}\newline{\footnotesize\color{muted}\file{config.js}} & \code{60_000} & Daugiausia teksto, kurį gali nešti viena naudotojo žinutė. \\
\code{MAX_PARTS}\newline{\footnotesize\color{muted}\file{routes/chat.js}} & \code{30} & Daugiausia dalių, kurias gali turėti viena naudotojo žinutė. \\
\code{MAX_REQUEST_BYTES}\newline{\footnotesize\color{muted}\file{config.js}} & \code{"512kb"} & Didžiausias viešos pokalbio užklausos turinys. \\
\code{MAX_HISTORY_MESSAGES}\newline{\footnotesize\color{muted}\file{config.js}} & \code{200} & Ilgiausia pokalbio šaka, atkuriama iš API; senesnės žinutės atmetamos. \\
\code{MAX_HISTORY_CHARS}\newline{\footnotesize\color{muted}\file{config.js}} & \code{400_000} & Kiek žinučių teksto vienas ėjimas gali nusiųsti modeliui (apie 100\,000 žetonų); pirmiausia atmetami ištisi seni ėjimai. \\
\code{MAX_STORED_CHARS}\newline{\footnotesize\color{muted}\file{services/transcript.js}} & \code{380_000} & Prieš išsaugant atsakymas sutraukiamas žemiau šios ribos (API atmeta eilutę, ilgesnę nei 400\,000). \\
\code{RATE_LIMIT_WINDOW_MS}\newline{\footnotesize\color{muted}\file{config.js}} & \code{60_000} & \code{RATE_LIMIT_TURNS} langas. \\
\code{TITLE_TIMEOUT_MS}\newline{\footnotesize\color{muted}\file{config.js}} & \code{8_000} & Ilgiausias laikas, kiek gali trukti pokalbiui pavadinimą suteikiantis iškvietimas; jį viršijus lieka automatinis pavadinimas. \\
\code{TITLE_INPUT_CHARS}\newline{\footnotesize\color{muted}\file{config.js}} & \code{2_000} & Kiek pirmosios žinutės simbolių perskaito pavadinimo modelis. \\
\code{TITLE_MAX_CHARS}\newline{\footnotesize\color{muted}\file{services/title.js}} & \code{80} & Ilgiausias išlaikomas pavadinimas. \\
\code{MAX_LISTED}\newline{\footnotesize\color{muted}\file{llm/prompt.js}} & \code{30} & Kiek leidimų, standartų, dokumentų ir rolių išvardija įmonės kontekstas; kai jų daugiau, modelis kviečia įrankį. \\
\code{MAX_REPORTS}\newline{\footnotesize\color{muted}\file{llm/prompt.js}} & \code{5} & Kiek įmonės ataskaitų išvardija kontekstas, naujausios pirmos. \\
\code{MAX_ACTIVITIES}\newline{\footnotesize\color{muted}\file{llm/prompt.js}} & \code{1_500} & Įmonės veiklos aprašymo simbolių skaičius kontekste. \\
\code{MAX_JSON_FIELD}\newline{\footnotesize\color{muted}\file{llm/prompt.js}} & \code{2_000} & Struktūrizuoto profilio lauko simbolių skaičius kontekste. \\
\code{MAX_ROLE_CHARS}\newline{\footnotesize\color{muted}\file{llm/prompt.js}} & \code{6_000} & Į instrukciją įdedamo rolės profilio simbolių skaičius. \\
\code{CACHE_MS}, \code{RETRY_MS}\newline{\footnotesize\color{muted}\file{services/prompts.js}} & \code{60_000}, \code{5_000} & Kiek laiko pokalbis pakartotinai naudoja aktyvią instrukcijos versiją ir kaip greitai vėl bandoma po nepavykusio atnaujinimo: naujai aktyvuota pokalbio instrukcijos versija pradedama naudoti maždaug per minutę. \\
\code{MAX_TOOLS}\newline{\footnotesize\color{muted}\file{llm/mcp.js}} & \code{40} & Kiek įrankių vienas išorinis (MCP) serveris gali pridėti prie ėjimo. \\
\code{MAX_RESULT_CHARS}\newline{\footnotesize\color{muted}\file{llm/mcp.js}} & \code{20_000} & Kiek išorinio įrankio atsakymo simbolių gauna modelis. \\
\code{CALL_TIMEOUT_MS}, \code{CONNECT_TIMEOUT_MS}\newline{\footnotesize\color{muted}\file{llm/mcp.js}} & \code{30_000}, \code{8_000} & Ilgiausias laikas, kiek gali trukti vienas išorinio įrankio iškvietimas ir prisijungimas prie serverio su jo įrankių sąrašo gavimu. \\
klaidos tekstas (įterpta)\newline{\footnotesize\color{muted}\file{llm/tools/shared.js} (\code{ToolFailure})} & \code{400} & Kiek nepavykusio įrankio iškvietimo priežasties simbolių perskaito modelis. \\
\code{TIMEOUT_MS}\newline{\footnotesize\color{muted}\file{services/django.js}} & \code{15_000} & Kiek laiko asistento paslauga laukia, kol API atsakys į vieną užklausą; jį viršijus įrankis nepavyksta su \code{DJANGO_UNREACHABLE}. \\
\code{JUDGE_TIMEOUT_MS}\newline{\footnotesize\color{muted}\file{llm/tools/corpus.js}} & \code{150_000} & Toks pat laukimas vien nuorodų patikrai: API prieš atsakant paveda modeliui kiekvieną teiginį palyginti su jo nuostata. \\
\code{EVIDENCE_WAIT_MS}\newline{\footnotesize\color{muted}\file{llm/tools/companies.js}} & \code{1_500} & Kiek laiko \code{addCredential} laukia, prieš dar kartą bandydamas susieti įrodantį dokumentą, kuris dar nepriklauso įmonei. \\
\end{xltabular}


\section{Kaip keičiama konstanta}\label{sec:config:change}

Konstanta keičiama kode ir įsigalioja, kai ją laikantis procesas kitą kartą paleidžiamas:
gamyboje --- kai atnaujinimo scenarijus perkuria atvaizdus, kūrimo režime API ir asistento
paslaugai --- iš karto, nes jos persikrauna išsaugojus failą, o užduočių vykdytojo ciklui ---
tik jį paleidus iš naujo (\cref{ch:ops}). Jau vykdoma užduotis išlaiko reikšmę, su kuria
pradėjo, todėl dauguma konstantų veikia nuo kito vykdymo. Kelios konstantos susietos su
instrukcija, kurioje nurodytas tas pats skaičius, su kita konstanta arba su testu, ir jas būtina
keisti kartu (\cref{tab:config:depends}); kitas galima keisti atskirai.

\begin{table}[H]
\caption{Konstantos, nuo kurių priklauso instrukcija, kita reikšmė ar testas}\label{tab:config:depends}
\begin{tabularx}{\textwidth}{@{}L{50mm}Y@{}}
\toprule
\hd{Konstanta} & \hd{Kas nuo jos priklauso} \\
\midrule
\code{MAX_TIERS} (24) & Instrukcija \prompt{corpus.sanctions}: \enquote{Ne daugiau kaip dvidešimt keturios pakopos}. \\
\code{REQUIREMENTS_PER_DOCUMENT} (20) & Instrukcija \prompt{company.internal_requirements}: \enquote{Ne daugiau kaip dvidešimt reikalavimų}. \\
\code{CITATION_FIELDS} (citata 200) & Instrukcija \prompt{report.inventory}: citatos \enquote{ne ilgesnės kaip 200 simbolių}. \\
\code{ISO_CLAUSES}, \code{GUIDELINE_PARAGRAPHS} & Nuorodų sąrašas instrukcijoje \prompt{programme.raci}. \\
\code{MIN_STEPS}, \code{MIN_SCENARIOS}, \code{MIN_CATEGORIES} & Laikomos mažesnės, nei prašo instrukcijos \prompt{programme.raci} ir \prompt{programme.risk_scenarios}, kad geras atsakymas niekada nepažeistų taisyklės. \\
\code{SCHEMA_VERSION} & Būtina padidinti su kiekvienu inventoriaus schemos pakeitimu; \file{tests/test_citation_audit.py} tikrina rakto formą. \\
\code{GENERATE_TIMEOUT_S} ir \code{GENERATE_TIMEOUT_MS} & Keičiamos kartu: API laukia 30~s ilgiau nei asistento paslauga. \\
\code{EMBED_DIMENSIONS} & Teisės aktų bazės ir įmonės dokumentų vektorių stulpelių plotis duomenų bazėje; pakeitimui reikia migracijos ir visų vektorių perskaičiavimo. \\
\code{MAX_SERIES_ARTICLES} (15) & \file{tests/test_citation_audit.py} tikrina, kad ji lygi 15. \\
Kiekviena \file{config.js} numatytoji reikšmė & Išvardyta ir patikrinta asistento paslaugos faile \file{tests/config.test.mjs}, kartu su kiekvieno skaitinio kintamojo minimumu. \\
\code{OBLIGATIONS_PER_CALL}, \code{DOCUMENT_PACK_CHARS}, \code{PLAN_FINDINGS_PER_CALL}, \code{CITATIONS_PER_CALL} & Konvejerio testai (\file{tests/test_pipeline.py}, \file{tests/test_assessment.py}, \file{tests/test_citation_audit.py}) jas nustato mažomis, kad priverstų susidaryti kelis paketus (\code{DOCUMENT_PACK_CHARS} atveju --- per didelį, kad būtų siunčiamas visas, dokumentų rinkinį); testai pateikiamų reikšmių nefiksuoja. \\
\bottomrule
\end{tabularx}
\end{table}

Konstanta, kuri keičia tai, ką gauna modelis, keičia ir tai, ką jis atsako, o testuose to niekas
nematuoja. Prieš tokį pakeitimą abu variantus palyginti padeda instrukcijų vertinimo rinkiniai
(\cref{ch:control}).
